\documentclass[11pt,a4paper]{article}
\usepackage[utf8]{inputenc}
\usepackage[T1]{fontenc}
\usepackage{amsmath,amssymb,amsthm}
\usepackage[margin=1in]{geometry}

\newtheorem{theorem}{Theorem}
\newtheorem{definition}{Definition}
\newtheorem{proposition}{Proposition}
\newtheorem{remark}{Remark}

\title{\textbf{S-wind: An Exact Algebraic Discovery System\\over $\mathbb{Q}$ with Closure Certification}}
\author{Zhuoli Ma (Jouly Mars)\\
\small Rogue Intelligence LNC.\\
\small \texttt{joulyman@gmail.com}}
\date{January 2026}

\begin{document}

\maketitle

\begin{abstract}
We present \textbf{S-wind}, a formal exploration system for mathematics and 
theoretical physics that operates natively over the field of rational numbers 
$\mathbb{Q}$. Unlike computation frameworks that optimize for speed, S-wind 
optimizes for \emph{verifiability}---producing auditable derivation traces 
rather than numeric outputs. All arithmetic uses \textbf{Extended Precision 
Rational (EPR)} numbers with mathematically unbounded precision.

Key contributions: (1) \textbf{EPR arithmetic} ensuring exact computation 
with zero floating-point error; (2) \textbf{Bounded closure certification} 
providing provable exhaustion of search space; (3) \textbf{Self-loop discovery} 
extending short signals to increase degrees of freedom; (4) \textbf{183 OEIS 
sequences} for known pattern recognition.

The system has been applied to explore fundamental physics hierarchies, 
discovering algebraic relations connecting the mass hierarchy, coupling 
hierarchy, and cosmological constant through a unified recurrence structure.
\end{abstract}

\section{Introduction}

\subsection{The Floating-Point Problem}

Standard numerical computation uses IEEE 754 floating-point arithmetic, which 
introduces:
\begin{itemize}
    \item Rounding errors at every operation
    \item Error accumulation over long computations
    \item Platform-dependent results
\end{itemize}

For fundamental mathematics and physics where \emph{exactness} is required, 
these limitations are unacceptable.

\subsection{The Discovery Definition Problem}

Existing pattern recognition systems can fit formulas to data, but lack:
\begin{itemize}
    \item A formal criterion for what constitutes a ``discovery''
    \item Comparison against known authoritative sources
    \item Proof of search exhaustion (vs. ``we tried many things'')
\end{itemize}

\subsection{Our Contribution}

S-wind addresses both problems through:
\begin{enumerate}
    \item \textbf{EPR Arithmetic}: Extended Precision Rational numbers with 
    arbitrary precision over $\mathbb{Q}$
    \item \textbf{Knowledge Base}: 183 OEIS sequences for pattern recognition
    \item \textbf{Bounded Closure}: Mathematical proof of search exhaustion
    \item \textbf{Self-Loop Exploration}: Automatic signal extension for DoF 
    improvement
\end{enumerate}

\section{Core Architecture}

\subsection{Extended Precision Rational (EPR)}

\begin{definition}[EPR]
An EPR number is represented as:
\begin{verbatim}
pub struct EPR {
    numerator: ExtendedInt,    // Arbitrary precision integer
    denominator: ExtendedInt,  // Always positive, coprime
}
\end{verbatim}
where \texttt{ExtendedInt} uses \texttt{Vec<u64>} for arbitrarily large integers.
\end{definition}

\textbf{Properties}:
\begin{itemize}
    \item All operations (add, subtract, multiply, divide) are \textbf{exact}
    \item No overflow: can represent $1000!$ precisely
    \item Automatic GCD reduction maintains canonical form
    \item Comparison is exact ($=$ not $|a-b| < \varepsilon$)
\end{itemize}

\subsection{Hypothesis Space $H$}

The search space is formally defined as:
\begin{equation}
H = \{\text{poly}_{\leq 10}, \text{recurrence}_{\leq 2}, \text{geometric}, 
\text{OEIS}_{183}\}
\end{equation}

\begin{itemize}
    \item $\text{poly}_{\leq 10}$: Polynomials up to degree 10
    \item $\text{recurrence}_{\leq 2}$: Linear recurrences of order $\leq 2$
    \item $\text{geometric}$: Geometric sequences and progressions
    \item $\text{OEIS}_{183}$: 183 fundamental sequences from OEIS
\end{itemize}

\subsection{Transformation Set $T$}

\begin{equation}
T = \{\Delta, \Delta^2, \Delta^3, \Sigma, \text{ratio}, |\cdot|, \text{sgn}, 
(\cdot)^2, n \cdot f, \text{MA}_3, \text{prod}, \text{mod}_{10}\}
\end{equation}

12 transformations enabling systematic exploration of signal structure.

\subsection{Bounded Closure Condition}

\begin{definition}[Closure]
Exploration terminates when bounded closure is achieved:
\begin{equation}
\text{CLOSURE}(E) \Leftrightarrow \forall s \in E, \forall \tau \in T: 
\tau(s) \in E \text{ or } |\tau(s)| < 2
\end{equation}
\end{definition}

This guarantees that no unexplored territory remains within the hypothesis 
space---a \textbf{proof of exhaustion}, not ``we gave up.''

\subsection{Discovery Criterion}

\begin{definition}[Discovery]
A result is a \textbf{discovery} if and only if:
\[
\text{Discovery} = (\text{Result} \neq \text{Known}) \land (\text{Verified}) 
\land (\text{Non-Trivial})
\]
where Non-Trivial requires:
\[
\text{DoF} = n - (d + 1) \geq 2
\]
\end{definition}

This prevents inflation of interpolation into ``discovery.''

\section{Fate Exploration Engine}

\subsection{Architecture}

\begin{verbatim}
         +---------------+
         |     Fate      |
         | (Unified Core)|
         +-------+-------+
                 |
    +------------+------------+
    |            |            |
+---v---+   +----v----+   +---v---+
|Explore|   |  Meta   |   | Self  |
|(Basic)|   |(Nested) |   |(Loop) |
+-------+   +---------+   +-------+
                 |
         +-------v-------+
         | KnowledgeBase |
         |  (183 OEIS)   |
         +---------------+
\end{verbatim}

\subsection{Self-Loop Discovery}

For short input signals (DoF $< 2$), Fate performs self-loop exploration:
\begin{enumerate}
    \item Detect pattern from initial signal
    \item Use detected formula to extrapolate new points
    \item Re-feed extended signal into exploration engine
    \item Repeat until DoF $\geq 2$ or pattern becomes unstable
\end{enumerate}

\textbf{Example}:
\begin{verbatim}
Input:  [3, 6, 9]        (3 points, DoF = 1, trivial)
Detect: f(n) = 3n
Extend: [3,6,9,12,...,33] (11 points)
Output: DoF = 9, non-trivial, EmergentLaw
\end{verbatim}

\section{Application: Physics Hierarchy Discovery}

S-wind was applied to explore fundamental physics hierarchies:

\subsection{Input Data}

\begin{itemize}
    \item $s_1 = 22$: $\log_{10}(M_P/m_e)$ (mass hierarchy)
    \item $s_2 = 42$: $\log_{10}(\alpha/\alpha_G)$ (coupling hierarchy)
    \item $s_3 = 122$: $\log_{10}(\Lambda_{\rm QFT}/\Lambda_{\rm obs})$ 
    (cosmological constant)
\end{itemize}

\subsection{Discovery}

S-wind exhaustive exploration discovered:
\begin{equation}
\boxed{s_{n+1} = 4s_n - 46}
\end{equation}

\textbf{EPR Verification}:
\begin{align}
4 \times 22 - 46 &= 88 - 46 = 42 \quad \checkmark \nonumber\\
4 \times 42 - 46 &= 168 - 46 = 122 \quad \checkmark
\end{align}

This relation connects three fundamental physics hierarchies through a single 
linear recurrence.

\subsection{Closure Certificate}

\begin{table}[h]
\centering
\begin{tabular}{|l|r|}
\hline
\textbf{Metric} & \textbf{Value} \\
\hline
Total iterations & 186 \\
Containers explored & 186 \\
Closure skips (duplicates) & 103 \\
Queue remaining at termination & 0 \\
Key discoveries & 20 \\
Computation time & 4.2 seconds \\
\hline
\end{tabular}
\caption{Exhaustive search statistics}
\end{table}

Queue = 0 at termination proves bounded closure was achieved.

\section{Comparison with Existing Systems}

\begin{table}[h]
\centering
\begin{tabular}{|l|c|c|c|c|}
\hline
\textbf{Feature} & \textbf{S-wind} & \textbf{Ramanujan M.} & \textbf{AI Feynman} & \textbf{PySR} \\
\hline
Exact arithmetic & \checkmark EPR & $\times$ float & $\times$ float & $\times$ float \\
Closure proof & \checkmark & $\times$ & $\times$ & $\times$ \\
Self-loop & \checkmark & $\times$ & $\times$ & $\times$ \\
OEIS integration & \checkmark 183 & $\times$ & $\times$ & $\times$ \\
Derivation trace & \checkmark & $\times$ & $\times$ & $\times$ \\
\hline
\end{tabular}
\caption{Comparison with existing discovery systems}
\end{table}

\section{Limitations}

\begin{enumerate}
    \item \textbf{Pattern Classes}: Currently supports polynomial, exponential, 
    reciprocal, factorial, recurrence. Continued fractions are future work.
    \item \textbf{Transcendentals}: $\pi, e, \sqrt{2}$ must remain symbolic.
    \item \textbf{OEIS Coverage}: 183 sequences, not complete OEIS database.
    \item \textbf{Cross-signal relations}: Current hypothesis space does not 
    include multiplication/division across different signals.
\end{enumerate}

\section{Conclusion}

S-wind demonstrates that \textbf{exact arithmetic} combined with 
\textbf{closure certification} enables rigorous algebraic discovery:
\begin{itemize}
    \item All computation is exact over $\mathbb{Q}$ (zero floating-point error)
    \item Search exhaustion is \emph{proven}, not assumed
    \item Self-loop mechanism enables discovery from minimal input
    \item Negative results (failures) are documented as evidence
\end{itemize}

The system has been applied to fundamental physics, discovering algebraic 
relations connecting mass hierarchy, coupling hierarchy, and cosmological 
constant through a unified recurrence structure.

\section*{Acknowledgments}

Parts of this manuscript were drafted with assistance from AI-based systems 
(large language models). The author takes full responsibility for the 
correctness of the results and interpretations.

\section*{Availability}

\begin{itemize}
    \item \textbf{Source}: \texttt{https://github.com/joulyman/s-wind}
    \item \textbf{License}: MIT
    \item \textbf{Language}: Rust
    \item \textbf{Run}: \texttt{cargo run --release --bin fate\_exhaustive\_latex}
\end{itemize}

\begin{thebibliography}{9}
\bibitem{ieee754} IEEE 754-2019. ``Standard for Floating-Point Arithmetic.''
\bibitem{ramanujan} Raayoni et al. (2021). ``Generating conjectures on 
fundamental constants with the Ramanujan Machine.'' \textit{Nature} 590, 67--73.
\bibitem{aifeynman} Udrescu \& Tegmark (2020). ``AI Feynman: A physics-inspired 
method for symbolic regression.'' \textit{Science Advances} 6(16).
\bibitem{oeis} OEIS Foundation. ``The Online Encyclopedia of Integer Sequences.'' 
https://oeis.org/
\end{thebibliography}

\end{document}
